\chapter{背景提升：绩点、英语、科研与竞赛}
\label{ch04}
\ist{背景提升}
\begin{keybox}
背景提升是保研里最容易被误解的一章。很多流传的说法——“六级不过不能保研”“必须有论文才有戏”“有国奖就能上岸”“专利受理就加分”——在制度层面都站不住脚：教育部层面的推免文件只要求“成绩优秀、有较强的创新意识和创新能力”，从未规定统一的四六级分数，也从未发布过全国通用的竞赛认定目录；认定权在学校和学院，标准写在各自的推免办法与实施细则里，校与校、院与院之间可以差出几个量级。

本章因此不给你一套通用答案，而是给你三样东西：一是\textbf{判断方法}——查本校推免办法（决定你有没有资格）、查目标院系当年接收通知（决定你能不能投）、查夏令营通知（通常要求最高）；二是\textbf{量级感}——学业成绩在多数学校占七成以上，而论文、竞赛、专利、荣誉这些专长加分的上限往往只有个位数分值，先补排名还是先堆成果，账要算清楚；三是\textbf{边界}——哪些材料会被审核答辩、哪些行为会被认定为学术不端。

读的时候请始终带着一个问题：这一条在我学校、我学院、我要申的项目里，具体是怎么规定的？本章引用的每一处分数、比例与目录，都标注了学校与年份，请把它当作检索线索，而不是结论本身。
\end{keybox}

\section{科研入门}
\begin{qa}{M0190}{大二进实验室晚不晚？怎么找愿意带本科生的老师}
\textbf{结论：不晚，大二是进实验室比较舒服的时点。} 大二下进组，到大三下正好有一年多的连续投入，能形成可写进简历的成果；太早（大一）常因基础课没学完而只能看，太晚（大三下）则来不及出结果。

找老师的可行路径：
\begin{enumerate}
  \item 先锁定方向——读本院教师近三年的论文标题，挑 3-5 位真正在做你想做的方向的老师；
  \item 邮件联系，正文控制在 300 字内：你是谁、上过哪些相关课、成绩排名、读过他哪篇论文、你想做什么、一周能投入多少时间；
  \item 附上成绩单与简历，不要群发明显模板化的信。
\end{enumerate}

\textbf{判断老师是否愿意带本科生，看公开信号}而不是听说：山东大学数学学院夏令营要求申请人填写“有意向交流的导师姓名（3 个以内）”，说明目标院校默认学生在主动联系导师；中国人民大学国学院甚至直接接受“课程论文、学年论文”作为代表性学术论文提交，说明本科生的早期成果是被认可的载体。
\nref{S0037}
\nref{S0006}
\end{qa}

\iin{山东大学}
\ith{科研入门}
\ist{背景提升}
\begin{qa}{M0191}{进组之后发现只是打杂，怎么才能有产出}
\textbf{结论：打杂的根因通常是“没有明确交付物”，不是老师不重视你。} 破局办法是主动把杂活收敛成一个可交付的东西。

三个可操作动作：
\begin{enumerate}
  \item 接手任何杂活时先问一句“这部分最后要交出什么形式的结果”——数据集、一段代码、一张图、还是文献综述；
  \item 把杂活升级为自己的小课题：整理文献就顺手写一篇综述初稿，跑数据就把方法写成可复现的脚本；
  \item 每两周给导师一段书面进度，写明你做了什么、卡在哪、下一步要什么支持。
\end{enumerate}

\textbf{学校对“参与”的认定是有具体要求的}。大连理工大学商学院在“研究潜力”项明确写道：若有多名参与者，需阐述申请者本人的工作及贡献；中南大学材料学院、安徽师范大学都要求专家审核小组排除“有名无实”，未通过审核鉴定或答辩的成果不得计入推免成绩。也就是说，挂名式的参与在材料审核环节就会被剔掉，只有你能说清自己那部分工作才算数。
\nref{S0549}
\nref{S0273}
\nref{S0646}
\end{qa}

\iin{大连理工大学}
\ith{科研入门}
\ist{背景提升}
\begin{qa}{M0211}{保研看重的科研潜力具体怎么体现}
\textbf{结论：科研潜力在评审表里被拆成了可打分的几块，不是玄学。} 你要做的是把抽象描述翻译成对应载体。

常见的四个载体：
\par\vspace{2pt}\nobreak\begin{center}\footnotesize
\begin{tabularx}{\textwidth}{>{\raggedright\arraybackslash}X>{\raggedright\arraybackslash}X}
\toprule
\textbf{载体} & \textbf{说明} \\ \midrule
论文 & 多数学校限定独立作者或第一作者，且要求期刊在指定目录内 \\
竞赛 & 需在本校认定的竞赛清单内、且达到规定等级 \\
专利 & 通常只认发明专利，并要求第一发明人与本校为专利权人 \\
科研训练 & 大创、课题参与，需能说清个人贡献 \\
\bottomrule
\end{tabularx}
\end{center}\vspace{2pt}\par

判断依据很直接：华东师范大学计算机学院夏令营把“有高质量学术论文公开发表、丰富的科研项目经历或在全国重大竞赛中获奖”列为科研能力的说明项；中南大学材料学院把科研能力直接量化成 JCR 分区论文、发明专利、竞赛名次的分值表，且论文多篇不累计、只取最高。教育部《关于加强硕士研究生招生复试工作的指导意见》也把“创新精神和创新能力”写进了专业素质考核范围。
\nref{S0108}
\nref{S0273}
\nref{S0046}
\end{qa}

\iin{华东师范大学}
\ith{科研入门}
\ist{背景提升}
\begin{qa}{M0212}{没有科研经历，面试时怎么回答相关问题}
\textbf{结论：没有正式科研经历时，不要硬造，改用“可验证的学术痕迹”。} 面试官真正想听的，是你有没有完整走过一次“提问—找方法—得出结论”的流程。

可以拿来替换的素材：
\begin{itemize}
  \item 课程设计或大作业：说清题目来源、你做了哪一部分、遇到的具体困难；
  \item 学年论文 / 课程论文：中国人民大学国学院在推免材料中明确接受“课程论文、学年论文”作为代表性学术论文提交，说明这类文本本身就是被官方认可的载体；
  \item 毕业设计选题与文献阅读量：能说出该方向三篇代表作并讲清差异；
  \item 自学课程的产出：代码仓库、读书笔记、复现报告。
\end{itemize}

被问到“你没有科研经历”时的回答框架：承认事实 → 说明你用课程项目补上了哪一步 → 给出你读过的具体文献和你的判断 → 表达想做的方向。教育部《关于加强硕士研究生招生复试工作的指导意见》把“大学阶段学习情况及成绩”列为专业素质考核的第一项，说明完整的学业表现本身就是评分项。
\nref{S0006}
\nref{S0046}
\end{qa}

\iin{中国人民大学}
\ith{科研入门}
\ist{背景提升}
\begin{qa}{M0213}{怎么把课程设计包装成科研经历}
\textbf{结论：可以改写，但改的是“表述与贡献归因”，不是事实。} 课程设计本来就是你做的工作，问题在于你写成了流水账而不是研究过程。

改写四步：
\begin{enumerate}
  \item 把题目改成一个问题：“基于 X 方法解决 Y 场景下的 Z 问题”；
  \item 补上背景与动机：为什么这个问题值得做，你参考了哪篇文献的方法；
  \item 明确你的个人贡献：用了什么工具、跑了多少数据、结果如何、指标是多少；
  \item 写清局限：你的方法在什么条件下失效——这一条最能体现科研思维。
\end{enumerate}

\textbf{边界很清楚}：大连理工大学商学院规定，参与人超过一人时“需阐述申请者本人的工作及贡献”；内蒙古大学、安徽师范大学、中南大学材料学院都要求对成果做审核鉴定并组织公开答辩，明确排除“有名无实”，未通过者不得计入推免成绩。把课程设计说成“科研项目”、把小组作业说成“独立完成”属于不实描述，是被明确堵死的方向。
\nref{S0549}
\nref{S0544}
\nref{S0646}
\nref{S0273}
\end{qa}

\iin{大连理工大学}
\ith{材料包装}
\ist{背景提升}
\begin{qa}{M0224}{文科没有实验室，科研经历可以怎么攒}
\textbf{结论：文科没有实验室，但科研载体比理工科更多，只是评价口径换成“文本与观点”。} 你要攒的是可提交的书面成果，而不是实验数据。

四类可攒的载体：
\begin{enumerate}
  \item \textbf{论文}：甘肃农业大学把 CSSCI、《中文核心期刊要目总览》收录期刊列入认可的学术期刊范围；安徽师范大学要求发表在“论文分类导向性期刊目录”C2 及以上期刊；大连理工大学商学院在科研论文项中同时列出中文核心（5 分/篇）与 EI 检索会议论文（3 分/篇）。
  \item \textbf{课程论文与学年论文}：中国人民大学国学院在推免申请材料中明确接受“课程论文、学年论文均可”。
  \item \textbf{学科竞赛}：文科类以案例大赛、调研报告、辩论与写作类赛事为主，但需落在本校认定的竞赛清单内。
  \item \textbf{社会调查}：田野调研、口述史、问卷与访谈的原始资料整理本身就是学术训练。
\end{enumerate}

操作上最容易见效的顺序是：先把一门专业课的期末论文写到能投的水平，再带着它去找老师改。
\nref{S0059}
\nref{S0646}
\nref{S0549}
\nref{S0006}
\end{qa}

\iin{甘肃农业大学}
\ith{科研入门}
\ist{背景提升}
\begin{qa}{M0230}{保研推荐信找科研导师写还是任课老师写}
\textbf{结论：优先找科研导师；但多数学校对推荐人的职称有硬性要求，先看规则再选人。} 两者的差别在于：科研导师能写具体的工作细节，任课老师通常只能写课程表现。

先看规则，因为很多项目直接限定推荐人资历：
\begin{itemize}
  \item 中国人民大学财政金融学院直博项目：需 2 位具有正高级及以上职称专家撰写推荐书；
  \item 北京大学信息管理系夏令营：需 2 位副教授（含）以上专家分别推荐；
  \item 教育部推免管理办法在规定特殊学术专长通道时，要求“三名以上本校本专业教授联名推荐”。
\end{itemize}

也就是说，\textbf{能写“细节”的人和符合“资历”的人有时不是同一个}。若两者不能兼得：科研导师写实质内容，另找一位职称够的老师补位。

推荐信的价值不在头衔，而在于对方能否说出你做过什么、做得多好。
\nref{S0018}
\nref{S0035}
\nref{S0454}
\end{qa}

\iin{中国人民大学}
\ith{推荐信}
\ist{背景提升}
\begin{qa}{M0231}{怎么在短时间内和老师建立能写推荐信的关系}
\textbf{结论：短时间能建立的不是“交情”，而是“可信的学术接触记录”。} 与其想办法请老师吃饭，不如让老师在一件事上真实地见过你的能力。

三条高效路径：
\begin{enumerate}
  \item \textbf{修他的课并且做出成绩}：课后带着具体问题去答疑、把课程论文写好再请他指导修改——这是成本最低的建立关系方式；
  \item \textbf{带着产出找他}：一封邮件里附上你写的文献综述或复现报告，请求点评，并说明你愿意投入的时长；
  \item \textbf{让他做一次真实的评价}：请他给你的报告提修改意见，来回两三轮，他就有资格写出细节。
\end{enumerate}

规则上的抓手：中国人民大学国学院在推免材料中要求“直博申请人须提供 2 份专家推荐书，推荐专家必须是在报考领域取得正高职称”，且在读材料中要求提交体现学术水平的代表性论文——也就是说，\textbf{让老师在具体材料上见过你，本身就是被制度鼓励的路径}。华北电力大学 2027 年接收办法规定，申请人须按报考学院的具体要求提交材料，各学院按学院（或学科专业）成立材料审核专家组对材料独立打分并择优确定复试名单——推荐信的具体要求分散在各学院通知里，动笔前先查目标学院的通知。
\nref{S0006}
\nref{S0415}
\end{qa}

\iin{中国人民大学}
\ith{推荐信}
\ist{背景提升}
\begin{qa}{M0233}{保研资格里综合素质怎么体现，需要准备什么材料}
\textbf{结论：综合素质由各校各院自定指标体系，但支撑材料高度趋同。} 你要准备的不是一堆证书，而是一份“可被逐条核对”的材料清单。

常见的四类支撑：
\par\vspace{2pt}\nobreak\begin{center}\footnotesize
\begin{tabularx}{\textwidth}{>{\raggedright\arraybackslash}X>{\raggedright\arraybackslash}X}
\toprule
\textbf{类别} & \textbf{典型材料} \\ \midrule
学业 & 成绩单、专业排名证明（须教务盖章） \\
科研 & 论文首页与目录、录用通知、项目结题证明 \\
荣誉 & 奖学金证书、荣誉称号文件 \\
服务 & 志愿服务证明、国际组织实习证明、参军入伍材料 \\
\bottomrule
\end{tabularx}
\end{center}\vspace{2pt}\par

依据：南开大学把“综合素质评价指标体系由各学院根据本学院实际情况设定”，可纳入特殊学术专长、科研训练表现、参军入伍、志愿服务、国际组织实习等因素；湖南大学把专项评价成绩的范围写为科研成果、省部级竞赛获奖、荣誉表彰、志愿服务、国际组织实习和参军入伍、双创三大赛、体育艺术类比赛国家级及以上奖励，且需答辩并全程录音录像；中南大学材料学院则按项目分别要求本科生院、校团委、学工部、体育教研部审核认定后公示。

准备时按“每一项都要有出具单位盖章”的标准来收集。
\nref{S0561}
\nref{S0540}
\nref{S0273}
\end{qa}

\iin{南开大学}
\ith{综测}
\ist{背景提升}
\begin{qa}{M0236}{保研目标校看重本科学校的课程难度吗}
\textbf{结论：目标校普遍不设置“课程难度系数”，但会通过其他方式对本科出身打折。} 换句话说，难度不被直接换算，但学校层次会被参考。

依据：
\begin{itemize}
  \item 教育部《推免工作管理办法（试行）》对推荐条件的表述是“勤奋学习，刻苦钻研，成绩优秀”，全文没有课程难度调整条款；
  \item 北京大学校级实施办法只说“原则上学业成绩处于下游的学生不予推荐”，综合排名标准由院系自行制定和解释，同样没有难度系数；
  \item 但接收端会有出身参考：山东大学经济研究院在夏令营通知中写明，若所在学校属于教育部公布的世界一流大学建设高校，或所在专业在教育部第四轮学科评估中结果为 B+ 及以上，成绩排名可适当放宽。
\end{itemize}

实际影响是：\textbf{课程难度低不会帮你加分，但你的排名会被放到本科院校的背景里一起看}。
\nref{S0454}
\nref{S0554}
\nref{S0026}\par\textbf{本版实操补充：} 课程难度通常不是可单独换算的统一加分项，招生单位更容易核对成绩单、课程内容和你对基础知识的掌握。若想证明“低分不是因为能力差”，准备三类可核验材料：

\begin{itemize}
  \item 课程大纲/培养方案：标出课程性质、学分、先修要求、考核方式，说明课程覆盖的知识点；
  \item 排名证明：同一专业/同一课程群里的相对位置，比单科绝对分更能说明问题；
  \item 作品/项目/课程设计：把课程训练转成可展示的产出，面试时讲清你负责的部分和决策过程。
\end{itemize}

不要只用“本校给分严”解释低分，也不要把课程难度当成万能借口；如果目标项目公开了先修要求，按清单补课或用MOOC/项目证明掌握程度。
\end{qa}

\iin{教育部学生服务与素质发展中心（学信网/研招网）}
\ith{院校偏好}
\ist{背景提升}
\begin{qa}{M0237}{自学了目标方向的课程，怎么在申请里体现}
\textbf{结论：自学课程要落在“可核验的凭证”上，不能只写“自学了某某课”。} 评审看不到你的学习过程，只看得到凭证。

三种可用的凭证：
\begin{enumerate}
  \item \textbf{成绩单上的记录}：教育部《普通高等学校学生管理规定》明确，学生可以申请辅修校内其他专业或选修其他专业课程，也可以参加学校认可的开放式网络课程学习，经学校审核同意后修读的课程成绩（学分）予以承认。能进成绩单的课，就是最有力的证明。
  \item \textbf{辅修/课程证明}：山东大学经济研究院在报名材料清单中直接列出“跨专业学生辅修经济学类课程证明”，说明辅修证明是接收方预期的材料之一。
  \item \textbf{产出型证明}：把自学内容做成可展示的东西——代码仓库、复现报告、文献综述、公开笔记，并在个人陈述里写清你用它解决了什么问题。
\end{enumerate}

写材料时的原则：先给凭证，再给结论；不要用“熟练掌握”这类无法验证的措辞。
\nref{S0040}
\nref{S0026}
\end{qa}

\iin{中华人民共和国教育部}
\ith{材料包装}
\ist{背景提升}
\section{社工}
\begin{qa}{M0206}{实习对保研有帮助吗？还是只对就业有帮助}
\textbf{结论：普通企业实习对推免资格基本不加分，但对面试与部分项目的材料有帮助；有明文加分的是“国际组织实习”。} 把两者分开看，才不会押错。

依据：
\begin{itemize}
  \item 各校推免加分条款中，“实习”几乎只以\textbf{国际组织实习}的形式出现。内蒙古大学把“到国际组织实习（三个月及以上）”列为同等条件下优先考虑的对象；安徽师范大学把“根据教育部公布的国际组织实习岗位，到国际组织实习三个月及以上”单列为国际交流类专长分；大连理工大学商学院在“国际视野”项中规定，到代表性国际组织实习半年以上计 5 分，公派国际交流长期 5 分、短期 2 分。
  \item 普通企业实习未见推免加分条款，但教育部《普通高等学校学生管理规定》第十七条规定，参加社会实践等活动的经历、成果可以折算为学分，具体办法由学校规定。
\end{itemize}

所以实习的价值主要在两方面：给你的申请材料提供真实的项目经历，以及帮你想清楚要不要读研。
\nref{S0544}
\nref{S0646}
\nref{S0549}
\nref{S0040}
\end{qa}

\iin{内蒙古大学}
\ith{实习}
\ist{背景提升}
\begin{qa}{M0207}{去大厂实习会不会耽误保研准备时间}
\textbf{结论：风险不在“耽误时间”，而在挤占推免流程的关键窗口和必修课。} 大厂实习本身不扣分，但它会和排名的计分时点撞车。

两个具体风险：
\begin{enumerate}
  \item \textbf{计分时点}：湖南大学2027年的日程是 7 月 21 日前完成推免资格审核与学业综合成绩排名公示、8 月 28 日前学生申请与专项评价公示、8 月 31 日前推免遴选总成绩排名公示；中南大学材料学院要求成绩审核以截止 9 月 10 日 17:00 教务系统公布的加权平均分和专业排名为准。也就是整个大三暑假都在推免流程里。
  \item \textbf{必修课与首次成绩}：大连理工大学商学院要求修完前三学年培养方案所有必修课，且一考不及格的必修课不能超过 2 门；内蒙古大学要求所修培养方案内课程均已合格。实习期若出现课程不及格，是不可逆的损失。
\end{enumerate}

\textbf{交换/交流的成绩还有特殊规则}：大连理工大学商学院规定交流生在交流期间的原始成绩与转换成绩均不计入推免计算。实习虽不适用此条，但同理说明——不在本校修的课，成绩上的收益可能落空。
\nref{S0540}
\nref{S0273}
\nref{S0549}
\nref{S0544}
\end{qa}

\iin{湖南大学}
\ith{实习}
\ist{背景提升}
\begin{qa}{M0208}{社工经历在保研面试里会被问吗}
\textbf{结论：会被问，但通常不是主考项，除非它与你的专业方向有关。} 面试官看的是这段经历能不能反映你的组织力、执行力或问题解决能力。

规则依据：教育部《关于加强硕士研究生招生复试工作的指导意见》把综合素质和能力界定为“本学科（专业）以外的学习、科研、社会实践（学生工作、社团活动、志愿服务等）或实际工作表现等方面的情况”，并列出事业心、责任感、纪律性、协作性和心理健康情况。也就是说，社工经历属于被正式纳入考察范围的内容。

落地到面试环节，形式各不相同：大连理工大学商学院综合面试要求 5 分钟 PPT 介绍，内容须包含专业排名、本科课程及成绩、外语能力、参与科研创新类项目与竞赛情况等——在这一类结构化自我介绍里，社工往往只用一两句话带过。北京师范大学人工智能学院的面试则按专业素质与能力、综合素质与能力、英语口语听力分开打分。

准备方法：用一句话说清你的角色，再用一个具体事例说明你解决了什么冲突或推动了什么事，不要复述职位名称。
\nref{S0046}
\nref{S0550}
\nref{S0072}
\end{qa}

\iin{中华人民共和国教育部}
\ith{社工}
\ist{背景提升}
\begin{qa}{M0209}{当过班长团支书对保研有实质帮助吗}
\textbf{结论：有分数，但上限很低，且不能单独成为推免理由。} 班长、团支书这类岗位的收益主要通过“荣誉表彰”和“思想素养/综合素质”两个口子兑现。

规则依据：
\begin{itemize}
  \item 教育部《推免工作管理办法（试行）》第十四条第七项明确：在制定综合评价体系时，可对文艺、体育及社会工作特长等因素予以适当考虑，“但具备这些特长者必须参加综合排名，不得单列”。这决定了社工不可能成为绕过成绩的通道。
  \item 换算成具体分值更直观：中南大学材料学院荣誉称号加分中，校优秀学生干部 0.3 分、校优秀学生干部标兵 0.4 分，且多项荣誉不累计、只取最高；学院评价加分累计不超过 2 分。
  \item 大连理工大学商学院的“思想素养”满分 5 分，其中一部分由前三学年所获精神文明奖学金、社会工作奖学金、社会实践奖学金的数量换算，另一部分来自班级民主评议（1.2 分）、辅导员评议（0.6 分）、班主任评议（0.2 分）。
\end{itemize}

结论：班长团支书的实质帮助，是让你更容易拿到那几张荣誉证书，而不是直接给你保研名额。当与绩点冲突时，优先保绩点。
\nref{S0454}
\nref{S0273}
\nref{S0549}
\end{qa}

\iin{教育部学生服务与素质发展中心（学信网/研招网）}
\ith{社工}
\ist{背景提升}
\begin{qa}{M0210}{志愿活动支教经历算不算保研软实力}
\textbf{结论：分情况——志愿服务本身是“加分项”而非“软实力装饰”，支教则是一条独立的推免通道。} 二者性质完全不同，别混为一谈。

\textbf{志愿服务}：多校把它写进加分指标，但门槛不低。安徽师范大学要求“每学年坚持参加志愿服务活动并获得省级以上志愿服务先进个人表彰，或在国家级重大活动中表现突出”；中南大学材料学院荣誉称号加分中，校优秀青年志愿者 0.3 分、省优秀大学生社会实践个人 0.8 分，且不累计只取最高。

\textbf{研究生支教团}：这是占用教育部专项计划名额的独立推免途径，不是加分。西南交通大学第29届支教团的报名条件为：完成前三学年（五年制前四学年）必修课、综合测评成绩不低于 80 分、专业年级排名前 50\%、全国大学英语四级不低于 425 分，且报名时须已通过中小学教师资格证笔试全部科目。安徽工业大学支教团则要求学业成绩专业排名前 35\%。

注意：教育部《2026年全国硕士研究生招生工作管理规定》中“服务期满且考核合格可在初试总成绩加 10 分”的政策适用于\textbf{统考}，与推免不是同一条路径。
\nref{S0646}
\nref{S0273}
\nref{S0061}
\nref{S0589}
\nref{S0051}
\end{qa}

\iin{安徽师范大学}
\ith{社工}
\ist{背景提升}
\begin{qa}{M0227}{出国交换一学期对保研是加分还是减分}
\textbf{结论：加分有限，风险不低，关键看本校的学分认定规则。} 交换不是天然的加分项。

\textbf{风险面（更要紧）}：大连理工大学商学院明确规定，“交流生在交流期间原始成绩和转换后的成绩均不计入推免计算，含国内（包括港澳台地区）和国外高校”。安徽师范大学的规则是：参加过国（境）内、外校际交流项目的学生，“须通过课程替代、学分认定完成交流期间本专业必修课程才可参加成绩计算和排名”。这两条合起来的意思是——\textbf{交换学期的课若不能替成本专业必修课，你可能既拿不到绩点，还可能因必修课未完成而失去资格}。

\textbf{收益面}：大连理工大学商学院“国际视野”项可计入公派国际交流（长期 5 分、短期 2 分）、在学院认定的学科重要国际会议上口头汇报论文（5 分）。

所以做决定前必须先做一件事：拿着交换学期的课程表，去教务处逐门确认能否替代本专业必修课、成绩如何转换、是否计入推免排名。
\nref{S0549}
\nref{S0646}
\end{qa}

\iin{大连理工大学}
\ith{交换}
\ist{背景提升}
\begin{qa}{M0228}{参加暑期学校或暑期科研营对保研有用吗}
\textbf{结论：对推免资格不加分，对获取接收单位的录取机会有实际作用，但作用取决于你从中拿到了什么。} 暑期学校与暑期科研营属于“过程性投入”，不是成果本身。

从证据看，各校推免加分条款里没有“参加暑期学校”这一项；加分只认成果——论文、竞赛、专利、荣誉、志愿服务、国际组织实习。

那它到底有什么用：
\begin{itemize}
  \item \textbf{建立学术接触}：山东大学数学学院夏令营要求申请人填写“有意向交流的导师姓名（3 个以内）”，说明入营本身就是一次导师匹配的机会；
  \item \textbf{拿到证明材料}：出访/交流类经历可以写进个人陈述，国际交流类还可能进入部分学校的国际视野评价项（大连理工大学商学院对公派国际交流长期计 5 分、短期 2 分）。
\end{itemize}

判断标准：一个暑期营值不值得去，看它结束后你能带走什么——一篇能投的论文初稿、一封导师推荐信、还是一段能讲满三分钟的科研细节。只有“参加过”这三个字，收益极低。
\nref{S0037}
\nref{S0549}
\end{qa}

\iin{山东大学}
\ith{暑研}
\ist{背景提升}
\begin{qa}{M0229}{海外暑研对保研申请有帮助吗}
\textbf{结论：对“申请”有帮助，对“推免资格”不加分。} 海外暑研的价值在于产出可写进材料的成果和一段可信的导师关系。

\textbf{加分侧}：现有院校条款中没有“海外暑研”单项，能兑现的只有其衍生成果。大连理工大学商学院“国际视野”项承认公派国际交流（长期 5 分、短期 2 分）和在认定的学科重要国际会议上口头汇报论文（5 分）；若目标是 C9 与中科院系统，英语门槛往往更高，海外经历需要配合英语成绩才能发挥作用。

\textbf{门槛侧}：北京大学计算机学院的夏令营申请条件之一就是英语水平达到四级 525 分、六级 450 分、托福 100 分、GRE 320 分或雅思 6.5 分等任一条，或作为第一（非共同一作）／第二作者（导师第一作者）发表过英文学术论文。可见顶尖院系更看重可验证的英语与论文，而不是“去过海外”。

\textbf{性价比判断}：如果这段暑研能产出一篇一作论文或一封有力推荐信，值得去；如果只是参观式体验，不如留在本校做完整课题。
\nref{S0549}
\nref{S0114}
\end{qa}

\iin{大连理工大学}
\ith{暑研}
\ist{背景提升}
\begin{qa}{M0234}{拿过国家奖学金对保研有没有实质加成}
\textbf{结论：取决于你所在学院是否把它写进加分条目；即使写进去，分值也很小且通常“不累计、取最高”。} 国奖是一份强荣誉，但不是推免的通行证。

看看它在评价体系里的实际位置：中南大学材料学院的荣誉称号加分表里，校级荣誉从 0.1 分到 0.6 分不等、“有多项荣誉获得者，不累计加分，取最高分”，且学院评价加分\textbf{累计不超过 2 分}；同一学院规定专项评价成绩各类型加分\textbf{累计不超过 10 分}。也就是说，荣誉类事项的整体天花板是个位数量级的分数，而学业成绩在多数学校占绝对主导（详见绩点一节的比例数据）。

真正有用的是它的\textbf{间接价值}：国家奖学金的评选本身就以学业成绩与综合表现排序为依据，能拿到它说明你的排名已经站在前列——而排名才是推免的硬通货；同时它在面试与材料审核中是一个可信度很高的能力信号。

结论：把国奖当成“排名靠前的结果”，而不是“用来换名额的工具”，会更接近实际。
\nref{S0273}
\end{qa}

\iin{中南大学}
\ith{荣誉}
\ist{背景提升}
\section{时间分配}
\begin{qa}{M0204}{大三了还没竞赛没科研，怎么补救}
\textbf{结论：先补门槛，再补专长；且要认清专长的天花板。} 大三才开始不等于没救，但要选对发力点。

\textbf{第一优先：确认自己是否还站在门槛内。} 各校的门槛是排名口径而非绝对分数：成都大学商学院要求本专业必修课程加权平均分排名前 20\%，安徽师范大学要求推免课程加权平均成绩排名前 35\%、且前三年综合测评排名前 35\%，内蒙古大学要求平均学分绩点排名不低于本专业前 50\%，并规定非英语专业四级 425 分。

\textbf{第二优先：用剩余时间集中做一件能形成成果的事。} 竞赛和大创的窗口在大三仍有：湖南大学 2027 年推免流程到 8 月底才完成专项评价公示，各校学科竞赛报名通常贯穿全年。

\textbf{专长的天花板要清楚}：内蒙古大学的论文与竞赛加分属于学分绩点加分、最高不超过 0.15；安徽师范大学专长成绩最高不超过 10 分；中南大学材料学院专项评价各类型累计不超过 10 分。所以补专长能拉开的差距有限，补排名和外语言门槛才是主线。
\nref{S0079}
\nref{S0646}
\nref{S0544}
\nref{S0540}
\nref{S0454}
\nref{S0273}
\end{qa}

\iin{成都大学商学院}
\ith{补救策略}
\ist{背景提升}
\begin{qa}{M0205}{保研前的暑假该做科研还是实习}
\textbf{结论：如果你缺的是“可被认定的成果”，选科研；如果你缺的是“经历与方向感”，实习的边际价值更取决于它是否属于国际组织实习。} 从加分口径看，两者并不对称。

\par\vspace{2pt}\nobreak\begin{center}\footnotesize
\begin{tabularx}{\textwidth}{>{\raggedright\arraybackslash}X>{\raggedright\arraybackslash}X>{\raggedright\arraybackslash}X}
\toprule
\textbf{选项} & \textbf{加分依据} & \textbf{分值/待遇} \\ \midrule
国际组织实习 & 内蒙古大学、安徽师范大学、中南大学材料学院均列入 & 3 个月以上可优先；安徽师大计 100 分专长分；大工商学院半年以上 5 分 \\
普通科研 & 论文、专利需达到期刊目录与署名要求 & 大工商学院 TOP 类 30 分/篇，中文核心 5 分/篇 \\
竞赛 & 需在本校认定清单内且达规定等级 & 大工商学院国赛一等奖 10 分 × 类别系数 × 排名系数 \\
\bottomrule
\end{tabularx}
\end{center}\vspace{2pt}\par

需要留意的时点：湖南大学 2027 年 8 月 28 日前完成学生申请与专项评价公示、8 月 31 日前完成推免遴选总成绩排名公示——整个暑假都在流程内，所以这个假期的安排应当服务于“在这两个时点前拿到可提交的结果”。一篇已被录用的论文、一个已公布的竞赛奖项，比一段还在进行中的实习或科研更有用。
\nref{S0549}
\nref{S0646}
\nref{S0544}
\nref{S0273}
\nref{S0540}
\end{qa}

\iin{大连理工大学}
\ith{时间分配}
\ist{背景提升}
\begin{qa}{M0214}{时间怎么分配？绩点科研竞赛英语社工全都要吗}
\textbf{结论：权重决定优先级——学业成绩是主干，其余全是枝叶。} 先看各校对权重的硬性规定：

\par\vspace{2pt}\nobreak\begin{center}\footnotesize
\begin{tabularx}{\textwidth}{>{\raggedright\arraybackslash}X>{\raggedright\arraybackslash}X>{\raggedright\arraybackslash}X}
\toprule
\textbf{学校} & \textbf{学业成绩占比} & \textbf{其余} \\ \midrule
甘肃农业大学 & 课程平均成绩占比不得低于 70\% & 奖励成绩（加分）不得高于 30\% \\
大连理工大学商学院 & 学习成绩 85\% & 思想素养 5\% + 综合素质与能力 10\% \\
成都大学商学院 & 学业成绩 90\% & 综合能力 10\% \\
南开大学 & 计算学分绩权重不得低于 50\% & 综合素质评价成绩由各学院设定 \\
安徽师范大学 & 推免课程加权平均成绩 + 专长成绩×0.1 & 专长成绩最高不超过 10 分 \\
\bottomrule
\end{tabularx}
\end{center}\vspace{2pt}\par

\textbf{加分的天花板普遍很低}：内蒙古大学论文/竞赛加分最高不超过 0.15（学分绩点口径）；中南大学材料学院各类型加分累计不超过 10 分、其中学院评价加分不超过 2 分。

所以合理顺序是：1）守住排名与必修课无不及格（这是资格线）；2）拿到外语门槛；3）再用剩余精力做一到两件能形成成果的专长；4）社工与志愿作为低成本补充。

注意一个制度性约束：教育部《推免工作管理办法（试行）》规定，文艺、体育及社会工作特长者“必须参加综合排名，不得单列”——不存在靠特长绕开成绩的路径。
\nref{S0059}
\nref{S0549}
\nref{S0079}
\nref{S0561}
\nref{S0646}
\nref{S0544}
\nref{S0273}
\nref{S0454}
\end{qa}

\iin{甘肃农业大学}
\ith{时间分配}
\ist{背景提升}
\begin{qa}{M0215}{大三还要不要继续做学生工作，会不会太占时间}
\textbf{结论：大三应减量而非直接退出，且优先保留能产出“荣誉凭证”的岗位。} 判断依据是分值，不是喜好。

社工类经历在评价体系里主要落在这几处：
\begin{itemize}
  \item 大连理工大学商学院“思想素养”满分 5 分，由前三学年所获精神文明奖学金、社会工作奖学金、社会实践奖学金的数量及思想品德评价构成，其中班级民主评议 1.2 分、辅导员评议 0.6 分、班主任评议 0.2 分；
  \item 中南大学材料学院荣誉称号加分中，校优秀学生干部 0.3 分、校优秀学生干部标兵 0.4 分、校优秀团干标兵 0.4 分，且多项荣誉不累计、取最高分；
  \item 教育部《推免工作管理办法（试行）》第十四条第七项要求，具备文艺、体育及社会工作特长者“必须参加综合排名，不得单列”。
\end{itemize}

对比一下量级：上述分值是 0.3—5 分这个量级；而学业成绩在甘肃农业大学占比不得低于 70\%、大连理工大学商学院占 85\%、成都大学商学院占 90\%。\textbf{时间有限时，把学生工作的投入压到“不卸任、能拿到荣誉”的最低成本，把释放出来的时间给绩点。}

一个例外情形：如果你走的是研究生支教团或“担任辅导员”类专项通道，学生工作经历的价值会显著上升。
\nref{S0549}
\nref{S0273}
\nref{S0454}
\end{qa}

\iin{大连理工大学}
\ith{时间分配}
\ist{背景提升}
\begin{qa}{M0216}{保研准备和专业课冲突，要不要翘课去跑夏令营}
\textbf{结论：缺课本身不直接扣分，但任何一门必修课不及格都可能直接出局。} 这是底线问题，不是性价比问题。

各校的硬约束：
\begin{itemize}
  \item 大连理工大学商学院：须修完前三学年培养方案要求的所有必修课程，且一考不及格的必修课不能超过 2 门，所有参评课程成绩均以一考成绩计算；
  \item 内蒙古大学：按照专业培养方案前三学年应修的必修课均已修完，所修培养方案内课程均已合格；
  \item 安徽师范大学：所有推免课程无补考或重修；
  \item 中南大学材料学院：推免生必须完成专业培养方案前 3 年所有教学环节且成绩合格，\textbf{有不及格课程者不得推荐}。
\end{itemize}

所以实际情况是：为了跑夏令营而缺席部分课程，风险不来自考勤，而来自你把期末考砸了。尤其是那些“只有期末一次成绩”的专业课，缺课积累的后果在期末集中兑现。

操作建议：跑夏令营前先确认该学期必修课的考核方式（是否有平时分、点名是否计入），把最不能丢的几门课的时间锁死，其余用补课与同学笔记补齐。
\nref{S0549}
\nref{S0544}
\nref{S0646}
\nref{S0273}
\end{qa}

\iin{大连理工大学}
\ith{时间分配}
\ist{背景提升}
\begin{qa}{M0217}{保研和考研能不能同时准备，怎么分配精力}
\textbf{结论：笔试层面可以并行，制度层面存在一个不可逆的分岔——一旦被录取为推免生，就不能再报名参加当年硕士统考。}

规则依据：
\begin{itemize}
  \item 教育部《2026年全国硕士研究生招生工作管理规定》与各校办法一致：已被招生单位录取的推免生不得再报名参加当年硕士研究生招生考试；
  \item 内蒙古大学实施办法第十五条写明“已被招生单位接收的推免生，不得再报名参加当年硕士研究生考试招生，否则取消其推免录取资格”；
  \item 西安交通大学2027年推荐免试研究生招生章程要求申请人在规定时间内完成“推免服务系统”的报名与录取流程，未按时完成视为放弃。
\end{itemize}

因此并行策略要按时间分段，而不是同时投入：
\begin{enumerate}
  \item \textbf{9 月前}：以推免为主线（夏令营、预推免、材料准备），统考复习只保留英语与数学的“保温”，每天固定一小段时间即可；
  \item \textbf{9 月底—10 月}：推免结果基本明确。此时要么已拿到接收函，专心收尾；要么进入统考全职冲刺，用两三个月补专业课与政治。
\end{enumerate}

判断依据很简单：推免窗口是刚性的、错过不再有；统考每年都有。9 月前时间应向推免倾斜。
\nref{S0051}
\nref{S0544}
\nref{S0548}
\end{qa}

\iin{中华人民共和国教育部}
\ith{时间分配}
\ist{背景提升}
\begin{qa}{M0221}{大三下还有必要刷绩点吗，还是全力准备材料}
\textbf{结论：先查清“计分截止时点”，再决定投多少时间给大三下的课。} 有的学校大三下成绩完全计入，有的只计入部分。

关键规则差异：
\begin{itemize}
  \item 中南大学材料学院：成绩审核\textbf{以截止 9 月 10 日下午 17:00 教务管理系统中公布的加权平均分和专业排名为准}，且补考和重修合格的课程成绩按 60 分计算；
  \item 大连理工大学商学院：\textbf{“截至公布综合成绩排名前一日，未出一考成绩的课程不纳入成绩计算范围”}——这意味着大三下未出分的课可能直接不计入，投入产出比要重算；
  \item 成都大学商学院：课程成绩计算采用\textbf{首次获得的成绩}，排名要求本专业必修课加权平均分前 20\%。
\end{itemize}

另外，推免流程本身会占用大三下到大四上的时间：湖南大学 2027 年 7 月 21 日前就要完成推免资格审核与学业成绩排名公示。

所以判断方法是：确认大三下课程在你们学院的计分范围与出分时间。如果计入，全力考；如果不计入，把时间切给材料与夏令营准备，但\textbf{保证不出现不及格}——那是不可逆的。
\nref{S0273}
\nref{S0549}
\nref{S0079}
\nref{S0540}
\end{qa}

\iin{中南大学}
\ith{时间分配}
\ist{背景提升}
\begin{qa}{M0241}{背景提升有没有优先级排序，先补哪块}
\textbf{结论：按“资格—主干—枝叶”三层排序，不要按“哪件事看起来更亮眼”排序。}

\textbf{第一层：守住资格线（不可逆，最高优先级）}
\begin{itemize}
  \item 必修课无不及格：中南大学材料学院明确“有不及格课程者不得推荐”，安徽师范大学要求所有推免课程无补考或重修；
  \item 排名进入入围比例：内蒙古大学不低于专业前 50\%、安徽师范大学前 35\%、成都大学商学院前 20\%；
  \item 外语门槛：内蒙古大学四级 425、安徽师范大学四级不低于 425（艺术体育类 355）。
\end{itemize}

\textbf{第二层：主干——学业成绩}
甘肃农业大学规定课程平均成绩占比不得低于 70\%，大连理工大学商学院为 85\%，成都大学商学院为 90\%，南开大学规定计算学分绩权重不得低于 50\%。这是唯一占比过半的项。

\textbf{第三层：枝叶——专长与加分}
上限普遍很低：内蒙古大学论文/竞赛加分最高不超过 0.15；安徽师范大学专长成绩不超过 10 分；中南大学材料学院各类型加分累计不超过 10 分、其中学院评价加分不超过 2 分。

正确的补课顺序是：\textbf{先把排名拉回安全区，再补外语，最后才堆成果}。反过来的顺序在多数学校拿不回排名上的差距。
\nref{S0273}
\nref{S0549}
\nref{S0646}
\nref{S0544}
\nref{S0079}
\nref{S0059}
\nref{S0561}
\end{qa}

\iin{中南大学}
\ith{时间分配}
\ist{背景提升}
\section{英语提升}
\begin{qa}{M0185}{四六级多少分才够保研用？425够吗}
\textbf{结论：没有全国统一的分数线，425 不是保研的通行证，它只是很多学校接收端的常见下限。}

先看全国层面：教育部《全国普通高等学校推荐优秀应届本科毕业生免试攻读硕士学位研究生工作管理办法（试行）》对推荐条件只写了“勤奋学习，刻苦钻研，成绩优秀……有较强的创新意识、创新能力和专业能力倾向”，\textbf{全文没有任何四六级分数规定}。所以“425 够不够”这个问题，只能问你要申请的学校。

再看校级条款，差异极大：

\par\vspace{2pt}\nobreak\begin{center}\footnotesize
\begin{tabularx}{\textwidth}{>{\raggedright\arraybackslash}X>{\raggedright\arraybackslash}X}
\toprule
\textbf{学校/项目} & \textbf{英语要求} \\ \midrule
内蒙古大学（推免资格） & 非英语专业四级 425 分；艺术体育类 355 分 \\
甘肃农业大学（推免资格） & 四级或六级 425 分 \\
成都大学商学院（推免资格） & 四级 425 或雅思 5.5 或托福 66 \\
中国人民大学财政金融学院 & 六级 426 分及以上 \\
北京大学计算机学院夏令营 & 四级 525 或六级 450，或托福 100、雅思 6.5 等任一条 \\
\bottomrule
\end{tabularx}
\end{center}\vspace{2pt}\par

\textbf{判断方法}：查三处——本校推免办法（决定你有没有资格）、目标院系当年接收通知（决定你能不能投）、以及夏令营通知（通常最严）。
\nref{S0454}
\nref{S0544}
\nref{S0059}
\nref{S0079}
\nref{S0018}
\nref{S0114}
\end{qa}

\iin{教育部学生服务与素质发展中心（学信网/研招网）}
\ith{英语门槛}
\ist{背景提升}
\begin{qa}{M0186}{六级没过能不能保研？哪些学校卡六级}
\textbf{结论：六级没过完全可以保研，甚至有些学校根本不设英语条款。} “六级不过不能保研”是一个被反复传播的错误说法。

\textbf{反例一：校级办法完全没有英语要求。} 北京大学《推荐优秀应届本科毕业生免试攻读研究生工作实施办法》列出的申请条件只有学籍、身心健康与成绩优良、无违纪记录、不计划赴境外留学几条，全文未提四六级，且明确“各院系的综合排名标准由院系自行制定和解释”。中山大学 2027 年接收推免办法也只把四级/六级成绩列为申请材料之一。

\textbf{反例二：要求的是四级，不是六级。} 内蒙古大学要求非英语专业四级 425 分；安徽师范大学要求非外国语学院学生四级不低于 425 分（艺术类、体育类 355 分）；西南交通大学研究生支教团要求四级不低于 425 分。

\textbf{确实卡六级的学校也存在}：中国人民大学财政金融学院要求六级 426 分及以上；南开大学法学院、马克思主义学院、环境科学与工程学院均要求六级 425 分及以上，\textbf{或}四级 550 分及以上。

所以正确的判断是：查你目标院系的通知原文，而不是听经验之谈。
\nref{S0554}
\nref{S0167}
\nref{S0544}
\nref{S0646}
\nref{S0061}
\nref{S0018}
\nref{S0133}
\end{qa}

\iin{北京大学}
\ith{英语门槛}
\ist{背景提升}
\begin{qa}{M0187}{保研要不要考雅思托福？多少分算有竞争力}
\textbf{结论：雅思托福不是必需项，但它是“替身条款”和“加分信号”两用——很多学校允许用它替代四六级过门槛。}

看几个明确的替代条款：
\begin{itemize}
  \item 中国人民大学财政金融学院：六级 426 分（含）以上，或具备同等水平的其他种类外语国家级考试成绩；
  \item 南开大学马克思主义学院等：六级 425 分，或托福 90 分、雅思 6 分及以上；
  \item 成都大学商学院：四级 425，或雅思 5.5，或托福 66；
  \item 浙江大学医学院直博：六级 480 分及以上，或新托福 90、雅思 6.0 及以上；
  \item 北京大学计算机学院夏令营：四级 525、六级 450、托福 100、雅思 6.5、GRE 320 等，符合其中一条即可。
\end{itemize}

\textbf{“多少分算有竞争力”没有统一答案}，取决于目标院系：多数学校的替代线落在雅思 6.0—6.5、托福 90—100 这个区间；越高层次的院系门槛越高。

\textbf{决策建议}：如果你四六级已经稳过线，雅思托福的边际收益主要体现在直博项目、英文授课项目和出国交换机会上；如果四六级始终卡线，考一次雅思往往比反复刷四六级更划算，因为多数学校接受替代。
\nref{S0018}
\nref{S0133}
\nref{S0079}
\nref{S0029}
\nref{S0114}
\end{qa}

\iin{中国人民大学}
\ith{英语提升}
\ist{背景提升}
\begin{qa}{M0188}{英语专业保研看专四专八吗？专八没过影响大吗}
\textbf{结论：取决于你申请的项目类型——外语类专业基本看专四专八，非外语类只把它当作四六级的替代品之一。}

\textbf{外语类/英语专业}：
\begin{itemize}
  \item 广东外语外贸大学规定，英语、商务英语、国际商务、外交学和教育学等专业的学生，全国英语专业四级考试成绩不低于 70 分，或托福不低于 95 分，或雅思不低于 7.0 分；
  \item 甘肃农业大学规定，英语专业学生全国高等学校英语专业四级或八级考试成绩达到 60 分及以上；
  \item 安徽师范大学规定，外国语学院学生须通过相应语种专业四级（或相当等级）。
\end{itemize}

\textbf{非外语类}：专四专八通常是英语门槛的平行选项，不是必须项。北京大学计算机学院夏令营的英语条件中，“国家英语专业四级或八级考试合格”与四级 525、六级 450、托福 100 等并列，符合其中一条即可；南开大学多个学院也把“英语专业四级（及以上）考试合格”列为可选条件之一。

\textbf{所以“专八没过影响大吗”要看身份}：英语专业学生，专八是核心指标；非英语专业学生，专八只是一个可选的英语水平证明，没过并不构成障碍，用六级或雅思托福替代即可。
\nref{S0134}
\nref{S0059}
\nref{S0646}
\nref{S0114}
\nref{S0133}
\end{qa}

\iin{广东外语外贸大学}
\ith{英语提升}
\ist{背景提升}
\begin{qa}{M0189}{保研面试有英文问答吗？口语差怎么办}
\textbf{结论：有，而且多数学校的推免复试明确把外语听说能力列为考核项，完全绕不开。} 但“口语差”不等于“注定失分”，因为口语大部分是结构化、可提前准备的。

\textbf{制度依据}：教育部《关于加强硕士研究生招生复试工作的指导意见》把“外语听说能力”列为专业素质和能力的四项考核内容之一；《2026年全国硕士研究生招生工作管理规定》第五十五条也写明外国语听力及口语测试在复试中进行，由招生单位自行组织、成绩计入复试总成绩。

\textbf{实际形式}（各校差异很大）：
\begin{itemize}
  \item 西安交通大学 2027 年推荐免试研究生招生章程：复试包含专业知识考核、\textbf{英语听力和口语测试}、思想政治素质和品德考核；
  \item 大连理工大学商学院：综合面试含英文自我介绍（3 分钟）；
  \item 北京师范大学人工智能学院：面试满分 100 分，其中英语口语和听力测试占 20 分；
  \item 南昌大学基础医学院：复试含外语听力和口语测试，总成绩相同时优先参考专业课成绩，其次为英语六级、四级成绩。
\end{itemize}

\textbf{准备方法}：英文自我介绍准备 1 分钟和 3 分钟两个版本；把常问的“为什么选这个方向”“介绍你的项目”写成英文稿并练到能自然说出；专业词汇按方向整理一页。
\nref{S0046}
\nref{S0051}
\nref{S0548}
\nref{S0550}
\nref{S0072}
\nref{S0001}
\end{qa}

\iin{中华人民共和国教育部}
\ith{英语提升}
\ist{背景提升}
\begin{qa}{M0225}{经管类保研要不要考CPA或CFA}
\textbf{结论：在检索到的推免办法、细则与接收通知中，未见任何院校把 CPA、CFA 列为推免加分项或资格条件。} 所以为保研去考 CPA/CFA，性价比需要重新评估。

\textbf{能查到的相关信息有限}：上海大学经济学院在学院介绍中提到自己是 CFA（注册金融分析师）官方合作伙伴，这是培养与课程层面的合作，并非推免加分条款；甘肃农业大学、内蒙古大学、安徽师范大学、中南大学材料学院等校的加分指标中，只有科研论文、竞赛获奖、发明专利、荣誉称号、参军入伍、志愿服务、国际组织实习几类，没有职业资格证书。

\textbf{那还要不要考？} 分场景：
\begin{itemize}
  \item 目标读学术型硕士、走科研路线：CPA/CFA 几乎不产生作用，同样的时间用来做一篇论文或一段助研更划算；
  \item 目标读专业型硕士（金融、审计、税务等）并计划就业：它在复试面试中可能是一个可聊的话题，在实习与求职中价值更直接；
  \item 时间极其充裕且不影响绩点：可以作为长期资产提前布局。
\end{itemize}

\textbf{结论}：把职业资格证当成“就业资产”而不是“保研筹码”来做决策，更接近实际。
\nref{S0327}
\nref{S0059}
\nref{S0544}
\nref{S0646}
\nref{S0273}
\end{qa}

\iin{上海大学经济学院}
\ith{证书}
\ist{背景提升}
\begin{qa}{M0232}{英语不好是不是注定保研吃亏}
\textbf{结论：不是。英语在推免里主要扮演“门槛”和“少数环节的权重项”，而不是总量第一名。}

\textbf{为什么不是硬伤}：
\begin{itemize}
  \item 一些学校的校级推免办法根本不设英语条款。北京大学 2017 年实施办法的申请条件只有学籍、学业成绩、品行等几条，没有四六级要求，综合排名标准由院系自定；中山大学 2027 年接收办法只把英语成绩列为申请材料之一。
  \item 在\textbf{资格认定}环节，英语通常是一道“过线即可”的闸门。内蒙古大学要求四级 425 分，安徽师范大学要求四级不低于 425 分（艺术体育类 355 分），达标即通过，不再比较分数高低。
\end{itemize}

\textbf{但它确实会在两处咬人}：
\begin{enumerate}
  \item \textbf{接收端门槛}：部分院系设六级门槛（如中国人民大学财政金融学院要求六级 426 分及以上），或设定很高的替代线（北京大学计算机学院夏令营要求托福 100 分或雅思 6.5 分等）；
  \item \textbf{复试环节}：多数学校把英语听力与口语测试计入复试成绩，北京师范大学人工智能学院面试中英语口语听力占 20 分。
\end{enumerate}

\textbf{所以真正的风险不是“英语差”，而是“英语拖到最后一刻才发现过不了目标院校的线”}。大三上学期前把四六级考到达标，风险就基本解除。
\nref{S0554}
\nref{S0167}
\nref{S0544}
\nref{S0646}
\nref{S0018}
\nref{S0114}
\nref{S0072}
\end{qa}

\iin{北京大学}
\ith{英语提升}
\ist{背景提升}
\begin{qa}{M0238}{保研前要不要考相关职业资格证}
\textbf{结论：分两类看——教师资格证对走研究生支教团/硕师计划几乎是必需，其他职业资格证未见推免加分条款。}

\textbf{有明文要求的场景}：
\begin{itemize}
  \item 研究生支教团：西南交通大学第 29 届支教团规定，报名时须已通过中小学教师资格证考试（自选一门学科）笔试部分全部考试，中共党员（含预备党员）或持有教师资格证者同等条件下优先；
  \item 硕师计划：安徽师范大学要求按时获得本科毕业证、学位证，\textbf{毕业时取得教师资格证}，师范类专业学生与参加过一学期实习支教的师范生同等条件下优先；
  \item 华中师范大学支教团选拔方式为综合素质面试，含说课、自我陈述与回答问题，实际考察的就是教学能力。
\end{itemize}

\textbf{没有依据的场景}：检索到的各校推免加分指标（科研论文、竞赛获奖、发明专利、荣誉表彰、参军入伍、志愿服务、国际组织实习）中，未见注册会计师、法律职业资格、计算机等级等证书单列。

\textbf{判断方法}：先确定你要走普通类还是专项通道。走普通类就按加分清单准备；走支教团/硕师计划，则教师资格证是硬性前置条件，且需提前一个考期准备。
\nref{S0061}
\nref{S0646}
\nref{S0614}
\end{qa}

\iin{西南交通大学}
\ith{证书}
\ist{背景提升}
\section{辅导班}
\begin{qa}{M0196}{专利对保研有用吗？实用新型和发明差别大吗}
\textbf{结论：专利有用，但“受理就加分”是明确错误的——“授权”才是通行口径，且多数学校只认发明专利并要求第一发明人。}

\textbf{各校的具体口径}：
\par\vspace{2pt}\nobreak\begin{center}\footnotesize
\begin{tabularx}{\textwidth}{>{\raggedright\arraybackslash}X>{\raggedright\arraybackslash}X}
\toprule
\textbf{学校} & \textbf{专利计分规则} \\ \midrule
大连理工大学商学院 & 发明专利 5 分/项、软件著作权 5 分/项、实用新型 3 分/项、外观型 1 分/项，最多限 1 项 \\
中南大学材料学院 & 以第一发明人获国家授权发明专利加 0.3 分，多项不累计、取最高 \\
成都大学商学院 & 以成都大学为专利权人、本人为第一发明人的\textbf{授权}发明专利方可纳入 \\
安徽师范大学 & 以第一发明人的身份获得国家\textbf{发明专利授权}，分值参照学校发明专利绩效认定表 \\
\bottomrule
\end{tabularx}
\end{center}\vspace{2pt}\par

\textbf{三点关键差异}：
\begin{enumerate}
  \item \textbf{发明 vs 实用新型差距很大}：大连理工给出的分值是 5 分对 3 分，且各校普遍把实用新型排除在外；
  \item \textbf{“受理/公开”不等于“授权”}：大连理工大学商学院明确“专利以授权日期为准”，成都大学、安徽师范大学均以“授权”为条件；
  \item \textbf{归属与位次是硬门槛}：普遍要求以本校为专利权人、本人为第一发明人，且成果须为在校期间取得。
\end{enumerate}

\textbf{判断方法}：查本校加分细则里的专利条目原文，看清楚它写的是“授权”还是“受理”、认不认实用新型。
\nref{S0549}
\nref{S0273}
\nref{S0079}
\nref{S0646}
\end{qa}

\iin{大连理工大学}
\ith{专利}
\ist{背景提升}
\begin{qa}{M0199}{大创项目对保研有什么用？结题了才有用吗}
\textbf{结论：大创本身在检索到的各校推免加分目录中未见单列，它的实际作用是“生产可被认定的成果”，而不是直接换成分数。}

\textbf{为什么这么说}：内蒙古大学、甘肃农业大学、安徽师范大学、中南大学材料学院、成都大学商学院的加分指标，认的是四类东西——科研论文、竞赛获奖、发明专利、荣誉表彰（含参军入伍、志愿服务、国际组织实习）。大创是\textbf{平台}，不是加分条目。

\textbf{所以关键在于大创产出了什么}：做出符合期刊目录要求且你是第一作者的论文，它按论文加分；作品参赛拿到本校竞赛清单内规定等级，它按竞赛加分；做出可授权的发明专利，它按专利加分。若只有一份结题报告，则在加分环节拿不到分，只能写进个人陈述作为科研经历。

\textbf{“结题了才有用吗”}：加分看的是成果是否达标（论文是否录用/发表、专利是否授权、竞赛是否获奖），而不是项目是否结题；但结题材料是证明“这段经历真实存在”的凭据。

\textbf{唯一的全国性衔接}：教育部《普通高等学校学生管理规定》第十七条规定，学生参加创新创业、社会实践等活动以及发表论文、获得专利授权等与专业学习、学业要求相关的经历、成果，可以折算为学分，具体办法由学校规定。
\nref{S0040}
\nref{S0544}
\nref{S0059}
\nref{S0646}
\nref{S0273}
\nref{S0079}
\end{qa}

\iin{中华人民共和国教育部}
\ith{大创}
\ist{背景提升}
\begin{qa}{M0218}{保研准备要不要报班？报班有用吗}
\textbf{结论：推免本身不需要报班，信息全部由官方渠道免费公开；需要付费购买的通常只是“整理好的信息”和“模板”。} 判断值不值，看它是否替你省下了本可以自己查的时间。

\textbf{先看清制度事实}：
\begin{itemize}
  \item 推免的名额分配、遴选办法、加分细则、公示名单都是法定公开内容。各校文件普遍要求办法与标准“公开透明”，内蒙古大学规定审定名单公示期不少于 10 个工作日，教育部《推免工作管理办法（试行）》要求推荐学校与接收单位均制订公开发布的标准与程序；
  \item 学校层面明确切割与商业机构的关系。东北师范大学在招生录取工作办法中声明：“我校不举办或参与举办任何形式的考研辅导活动，从未委托任何中介机构或个人开展招生信息发布、宣传、报名等相关工作。”
\end{itemize}

\textbf{所以报班的价值边界很清楚}：公开信息（各校接收通知、往年录取名单、加分目录）完全可以自己搜集；报班能提供的，主要是时间上的节省和文书修改意见。

\textbf{决策建议}：如果连本校推免办法都没读过，先别报班——读一遍会发现真正的问题比想象中少。若时间确实紧张，只买单项服务（文书润色、模拟面试）即可，不必上全套长周期课程。
\nref{S0454}
\nref{S0544}
\nref{S0607}
\end{qa}

\iin{教育部学生服务与素质发展中心（学信网/研招网）}
\ith{辅导班}
\ist{背景提升}
\begin{qa}{M0219}{保研文书找中介改值不值}
\textbf{结论：找中介“改”文书可以，“代写”不可以——这两件事在规则上是红线两侧。}

\textbf{红线在哪里}：
\begin{itemize}
  \item 各校对材料真实性的要求是刚性的。北京大学实施办法写明：有关被推荐者的学习成绩、综合能力介绍等材料要确保真实有据，\textbf{如有弄虚作假、改动数据者，一经发现查实，立即取消被推荐者的推免资格}；
  \item 材料审核有实质性的核验程序。中南大学材料学院、安徽师范大学、甘肃农业大学都要求成立专家审核小组，对科研成果、竞赛获奖等做审核鉴定并组织公开答辩，明确排除抄袭、造假、冒名及有名无实，未通过者不得计入推免成绩；
  \item 东北师范大学在招生办法中声明从未委托任何中介机构或个人开展招生信息发布、宣传、报名等工作，说明学校并不承认第三方代办的法律地位。
\end{itemize}

\textbf{所以真正有价值的中介服务只有一种：帮你把真实素材写清楚。} 也就是结构梳理、语言润色、篇幅压缩——内容必须是你自己的经历和数据。

\textbf{性价比判断}：如果你已经有完整素材，找人做一次结构梳理性价比尚可；如果你连素材都没有，中介无法凭空创造，硬写出来的东西在面试追问下会散架。
\nref{S0554}
\nref{S0273}
\nref{S0646}
\nref{S0059}
\nref{S0607}
\end{qa}

\iin{北京大学}
\ith{辅导班}
\ist{背景提升}
\begin{qa}{M0220}{保研要不要花钱买论文或竞赛奖，风险在哪}
\textbf{结论：买论文或买竞赛奖不仅风险极高，而且是被明文定性为学术不端的行为，一旦查实会直接取消推免资格并记入档案。}

\textbf{顶层规则（全国适用）}：
\begin{itemize}
  \item 教育部《高等学校预防与处理学术不端行为办法》（教育部令第 40 号）明确把“买卖论文、由他人代写或者为他人代写论文”列为学术不端行为，并要求高等学校对责任人作出处理；
  \item 教育部《普通高等学校学生管理规定》（教育部令第 41 号）规定，学位论文、公开发表的研究成果存在抄袭、篡改、伪造等学术不端行为情节严重的，或者代写论文、买卖论文的，学校可以给予开除学籍处分。
\end{itemize}

\textbf{推免环节的后果更直接}：
\begin{itemize}
  \item 甘肃农业大学规定，在推免过程中弄虚作假、存在学术不端行为的学生，一经查实取消推免资格、录取资格和学籍，并按规定记入《国家教育考试考生诚信档案》、通报所在学校、记入学生人事档案；
  \item 中南大学材料学院、内蒙古大学、成都大学商学院等均有同类条款；
  \item 而且很多学校在审核阶段就设置了实质核查：专家审核小组审核鉴定 + 公开答辩 + 全程录音录像 + 结果公示，未通过者不计分。
\end{itemize}

\textbf{真正的风险是}：用一次不可逆的代价，去交换一个分值上限通常只有个位数的加分项。
\nref{S0042}
\nref{S0040}
\nref{S0059}
\nref{S0273}
\nref{S0544}
\nref{S0079}
\end{qa}

\iin{中华人民共和国教育部}
\ith{学术风险}
\ist{背景提升}
\begin{qa}{M0226}{理工科要不要学编程或发专利来提升背景}
\textbf{结论：编程能力本身在推免加分目录中不单列，但对理工科面试（尤其是机试）有直接作用；专利可加分，但门槛高、周期长。}

\textbf{专利的分数账要算清楚}：
\begin{itemize}
  \item 大连理工大学商学院：发明专利 5 分/项、软件著作权 5 分/项、实用新型 3 分/项、外观型 1 分/项，最多限 1 项，且全部要求以第一发明人身份获得正式授权；
  \item 中南大学材料学院：以第一发明人获国家授权发明专利加 0.3 分，多项不累计取最高；
  \item 安徽师范大学：以第一发明人身份获得国家发明专利授权方可计入专长成绩。
  \item 共同点是：只认\textbf{授权}、要求\textbf{第一发明人}、普遍只认\textbf{发明专利}。实用新型在多数学校要么分值很低要么不被认可。
\end{itemize}

\textbf{编程能力的作用路径不同}：它不作为加分条目，但会直接影响考核环节。北京师范大学人工智能学院的复试就在面试之外单独设置上机测试（程序设计、数据结构与算法），综合考核成绩 = 机试成绩 + 面试成绩。对计算机与相关专业，机试失手是直接扣分，比任何加分项都硬。

\textbf{所以排序建议}：先确保你的核心专业课与机试能力不拖后腿，再考虑用专利去补专长加分。
\nref{S0549}
\nref{S0273}
\nref{S0646}
\nref{S0072}
\end{qa}

\iin{大连理工大学}
\ith{技能提升}
\ist{背景提升}
\begin{qa}{M0240}{内卷环境下怎么判断自己的背景在同档位里算什么水平}
\textbf{结论：用“同校同专业的公示名单”做对标，比任何经验判断都准。} 各校必须公开排名与加分数据，这是最硬的信息源。

\textbf{可直接查到的数据}：
\begin{itemize}
  \item 湖南大学计算机学院 2027 届推免名单公示包含学院、专业、学号、姓名、学业成绩、学业排名、综合成绩、综合排名、专业总人数、GPA 和备注，且逐条列出加分项与分值（如“第十六届蓝桥杯全国软件和信息技术专业人才大赛 省二等奖 +0.02”“湖南省大学生数学建模竞赛省级一等奖 排名第1 +0.05”）；
  \item 南京航空航天大学机电学院的拟推荐名单公示给出专业、学号、姓名与成绩；
  \item 湖南大学另行公示创新创业赛事、体育类、艺术类竞赛获国家级及以上奖励与志愿服务的加分名单及具体分值。
\end{itemize}

\textbf{对标方法}：
\begin{enumerate}
  \item 找到本校本专业近两年的推免公示，读你自己的排名位置与“最后一名被推荐者”之间的差距；
  \item 看加分项的分布——如果排在前面的人普遍只有 0.02—0.1 量级的加分，说明这个专业是“成绩决定论”，加分项无法逆转排名；
  \item 再去看目标院系的录取名单，看被录取者的本科院校与专业分布，判断自己的位置。
\end{enumerate}

\textbf{注意}：不要用“我拿了国奖”这类标签判断水平，要看它在你们专业加分细则里值多少分、能把你推进几名。
\nref{S0131}
\nref{S0100}
\nref{S0541}
\end{qa}

\iin{湖南大学}
\ith{定位方法}
\ist{背景提升}
\begin{qa}{M0242}{保研辅导机构说的背景提升项目靠谱吗}
\textbf{结论：区分两类服务——信息与文书辅导可以买，声称能“提供成果”的背景提升项目要高度警惕。}

\textbf{红旗信号一：承诺提供论文、专利、竞赛奖。} 教育部《高等学校预防与处理学术不端行为办法》把“买卖论文、由他人代写或者为他人代写论文”明确定性为学术不端；成都大学商学院、甘肃农业大学、安徽师范大学均要求专家审核小组对成果做审核鉴定并组织公开答辩，排除抄袭、造假、冒名及有名无实，未通过者不得计入推免成绩。\textbf{买来的成果在答辩环节是最先暴露的。}

\textbf{红旗信号二：打着“内部渠道”“合作院校”的名义。} 东北师范大学在招生录取工作办法中明确声明从未委托任何中介机构或个人开展招生信息发布、宣传、报名等相关工作；各校推荐名单、加分项都必须公示，公示期内接受异议与复议，不存在可以私下操作的通道。

\textbf{红旗信号三：用“保过”“兜底录取”当卖点。} 北京大学实施办法写明，材料如有弄虚作假、改动数据，一经查实立即取消被推荐者资格。

\textbf{可以买的部分}：信息整理、文书结构梳理与语言润色、模拟面试。这些是服务，不是成果。辨别方法很简单——看它交付的是“帮你表达”还是“替你拥有”。
\nref{S0042}
\nref{S0607}
\nref{S0554}
\nref{S0079}
\nref{S0059}
\nref{S0646}
\end{qa}

\iin{中华人民共和国教育部}
\ith{辅导班}
\ist{背景提升}
\section{竞赛认定}
\begin{qa}{M0197}{竞赛到底哪些在保研里被认？教育部的榜单在哪看}
\textbf{结论：不存在全国统一的“保研认可竞赛名单”，教育部层面的推免文件里没有竞赛目录。} 认定权在学校和学院手上，各校认定哪些赛事、认到什么等级，差别非常大。

\textbf{顶层文件确实没有目录}：教育部《全国普通高等学校推荐优秀应届本科毕业生免试攻读硕士学位研究生工作管理办法（试行）》全文只在推荐条件里写“有较强的创新意识、创新能力和专业能力倾向”，未列任何赛事名称；各年度全国硕士研究生招生工作管理规定同样没有推免竞赛目录。

\textbf{目录在各校自己的附件里}，形式也各不相同：
\begin{itemize}
  \item 大连理工大学商学院的竞赛计分注明“竞赛清单见附件”，并设 A 类（系数 1）、B 类（0.9）、C 类（0.7）的类别系数；
  \item 安徽师范大学限定为“学科竞赛分类导向性项目库（2023 版）”中奖励类别 B2 及以上的全国性、国际性比赛；
  \item 中南大学材料学院依据《中南大学推免生资格管理办法（试行）》（中大教字〔2021〕107 号）附件 1 所列“学科竞赛（含‘挑战杯’）获奖加分项目”；
  \item 成都大学商学院的口径是“国内外权威竞赛全国三等奖及以上奖励”。
\end{itemize}

\textbf{判断方法}：去本校教务处网站的推免专栏，找本校推免办法的附件——竞赛清单一定在那里，而不在教育部网站上。
\nref{S0454}
\nref{S0549}
\nref{S0646}
\nref{S0273}
\nref{S0079}
\end{qa}

\iin{教育部学生服务与素质发展中心（学信网/研招网）}
\ith{竞赛认定}
\ist{背景提升}
\begin{qa}{M0198}{数学建模国赛和美赛哪个更值钱}
\textbf{结论：无法一概而论，取决于本校竞赛清单是否列入美赛；但从加分口径看，国内赛事通常是基准，国际赛事“参照执行但不得低于国内要求”。}

\textbf{制度口径}：内蒙古大学实施办法写明，作为主力成员参加与学业相关的国内权威科研竞赛（全国赛）并获奖的，按一等奖不超过 0.15、二等奖不超过 0.12、三等奖不超过 0.1 的学分绩点加分，并注明\textbf{“国际赛事参照执行，但不得低于国内赛事相关要求”}。安徽师范大学的表述完全相同，且额外规定“全国赛不含省级赛或国家级赛事省级选拔赛和区域性比赛”。

这句话的实际含义是：\textbf{国际赛事不会自动比国内赛事更值钱，它要放到与国内赛事同一把尺子上量}。所以美赛能否加分、加多少，取决于你本校的清单把它归到哪个类别。

\textbf{可参考的量级差异}：湖南大学计算机学院 2027 届公示里出现的加分实例是省赛级别 0.01—0.05、国赛级别 0.1（如中国大学生计算机设计大赛国赛二等奖 +0.1）；中南大学材料学院对三大赛（中国国际大学生创新大赛、挑战杯、创青春）国家级金、银、铜奖分别至少加 0.2、0.15、0.1 分。

\textbf{判断方法}：拿本校清单对照两赛事的认定类别，再看往年公示名单里有没有人凭该项加过分。
\nref{S0544}
\nref{S0646}
\nref{S0131}
\nref{S0273}
\end{qa}

\iin{内蒙古大学}
\ith{竞赛认定}
\ist{背景提升}
\begin{qa}{M0200}{挑战杯互联网加拿省奖有用吗，还是必须要国奖}
\textbf{结论：省奖有没有用，取决于本校清单是否把省级赛事列入加分范围——在不少学校，省奖最多只是很低的加分，甚至完全不加分。}

\textbf{省奖被明确排除的情形}：安徽师范大学规定学术专长类的竞赛限“安徽师范大学学科竞赛分类导向性项目库（2023 版）”中奖励类别为 B2 及以上的全国性比赛和国际性比赛，并特别注明\textbf{“全国赛不含省级赛或国家级赛事省级选拔赛和区域性比赛”}。这条把省赛和分赛区赛都排除在外。

\textbf{省奖有加分的情形}：湖南大学计算机学院 2027 届公示中，湖南省大学生数学建模竞赛省级一等奖（排名第1）加 0.05、省二等奖加 0.02，而同一名单中国赛二等奖加 0.1，差距约 2—5 倍。中南大学材料学院规定由教育部等中央部委、省教育厅、省科技厅组织的重点赛事及学校重点赛事对应的省级比赛，一等奖排名第一加 0.5 分，二等奖 0.4 分，三等奖 0.3 分。

\textbf{所以答案是“有用，但量级小，且各校口径完全不同”}。湖南大学另一份专门公示显示，创新创业赛事（挑战杯）国家级奖励的加分明显更高：国家级特等奖排名第 1 加 6 分、国家级三等奖排名第 1 加 2 分，且随成员排名递减（第 3 名 2 分、第 4 名 1.5 分、第 5 名 1.2 分）。

\textbf{判断方法}：查本校清单的赛事级别定义，特别注意“省级赛”“分赛区”“选拔赛”是否被排除。
\nref{S0646}
\nref{S0131}
\nref{S0273}
\nref{S0541}
\end{qa}

\iin{安徽师范大学}
\ith{竞赛认定}
\ist{背景提升}
\begin{qa}{M0201}{哪些竞赛是水赛，参加了反而减分}
\textbf{结论：多校文件里确实存在“水赛”这个概念，并且有明确的排除机制——但判断标准不在赛事名字，而在本校清单与审核程序。}

\textbf{规则层面已经在防“水赛”}：
\begin{itemize}
  \item 安徽师范大学规定竞赛限“学科竞赛分类导向性项目库（2023 版）”中 B2 及以上的全国性、国际性比赛，且“全国赛不含省级赛或国家级赛事省级选拔赛和区域性比赛”；
  \item 中南大学材料学院直接写明：\textbf{“除了学校与学院重点赛事、排行榜重点赛事或教育部等中央部委、国家一级学会、省教育厅、省科技厅组织的赛事，其他赛事不予加分。”}
  \item 成都大学商学院只认定“国内外权威竞赛全国三等奖及以上奖励”；
  \item 内蒙古大学、广东外语外贸大学、安徽师范大学、云南师范大学都写入了同一条原则：\textbf{“对于社会质疑较多的赛事、刊物，各学院要从严把握。”}
\end{itemize}

\textbf{换句话说，“参加了反而减分”这个说法不准确——不达标的赛事是不加分，而不是扣分。} 真正的损失是你投入的时间和报名费。

\textbf{判断一个赛事是否值得打的三步}：1）是否在本校清单内；2）主办单位是否为教育部、中央部委、国家一级学会或省级主管部门；3）往年本校公示名单里有没有人凭它加分。三条有一条不满足，收益就要打折。
\nref{S0646}
\nref{S0273}
\nref{S0079}
\nref{S0544}
\nref{S0568}
\nref{S0134}
\end{qa}

\iin{安徽师范大学}
\ith{竞赛认定}
\ist{背景提升}
\begin{qa}{M0202}{竞赛要不要组队？队长和队员在保研认定上有区别吗}
\textbf{结论：有区别，而且各校都用明确的“排名系数”把差距写进了规则里——队长的加分通常是第 3、4 名成员的好几倍。}

看各校的具体算法：

\par\vspace{2pt}\nobreak\begin{center}\footnotesize
\begin{tabularx}{\textwidth}{>{\raggedright\arraybackslash}X>{\raggedright\arraybackslash}X}
\toprule
\textbf{学校} & \textbf{排名差异规则} \\ \midrule
湖南大学 & 挑战杯国家级特等奖按名次递减：第 1 名 6 分、第 3 名 2 分、第 4 名 1.5 分、第 5 名 1.2 分 \\
中南大学材料学院 & 排名第一按等级满分计，随成员排名依次递减 0.1 分 \\
大连理工大学商学院 & 作者排名系数前 6 名分别为 1、0.8、0.7、0.6、0.5、0.4 \\
安徽师范大学 & 团体只认定排名前 4 的成员，且按人数与位次比例分配分值 \\
\bottomrule
\end{tabularx}
\end{center}\vspace{2pt}\par

\textbf{关键结论有两条}：
\begin{enumerate}
  \item \textbf{位次决定分值}，队长与第 4、5 名的差距可以是 3—5 倍；
  \item \textbf{很多学校对成员人数设上限}，安徽师范大学团体只认定排名前 4 的成员，排在后面的成员在推免认定中等于没有贡献。
\end{enumerate}

\textbf{所以组队策略要反过来做}：先看本校细则规定的“认定人数上限”和“排名系数表”，再决定队伍规模。5 人以上团队里的后排成员，即使拿了国奖也可能无法计入加分。
\nref{S0541}
\nref{S0273}
\nref{S0549}
\nref{S0646}
\end{qa}

\iin{湖南大学}
\ith{竞赛认定}
\ist{背景提升}
\begin{qa}{M0203}{竞赛证书上没我名字能不能算}
\textbf{结论：证书上没有你的名字，基本不能计入加分；但团队型赛事有例外路径，前提是能出具学校认可的证明。}

\textbf{一般规则}：各校认定都要求可核验的署名。甘肃农业大学规定高水平科研创新成果应为第一作者发表，或与指导教师联合发表（限指导教师和学生分别为第一、二作者），且本校为第一署名单位；成都大学商学院要求以本校为参赛单位独立或作为第一完成人获得奖励；安徽师范大学规定团体只认定排名前 4 的成员。

\textbf{团队型赛事的例外}：中南大学材料学院对“未列出获奖者姓名的团体型赛事”设有专门条款——\textbf{取对应等级最高分值的一半，具体参与人以学校相关机构出具的证明和证书为准}。也就是说，这类情况需要学校相关部门出具参与证明，且分值要打对折。

\textbf{论文的类似条款}：中南大学材料学院规定，存在共同第一作者时，排名第一的按 80\% 计分，其他共同第一作者按 20\% 计分。

\textbf{底线提醒}：各校审核都会明确排除“冒名”与“有名无实”。成都大学商学院、内蒙古大学、安徽师范大学均要求专家审核小组审核鉴定并组织公开答辩，未通过者不得计入推免成绩。把别人的成果写进自己材料属于学术不端，风险远大于收益。
\nref{S0059}
\nref{S0079}
\nref{S0646}
\nref{S0273}
\end{qa}

\iin{甘肃农业大学}
\ith{竞赛认定}
\ist{背景提升}
\begin{qa}{M0239}{大三参加竞赛还来得及吗，会不会影响夏令营}
\textbf{结论：来得及，但要分清“为了加分”和“为了申请材料”两个目的，且要避开与夏令营流程的时间冲突。}

\textbf{时间上的硬约束}：推免流程本身占满大三暑假。湖南大学 2027 年的安排是 8 月 28 日前学生申请与专项评价公示、8 月 31 日前完成推免遴选总成绩排名公示，7 月 21 日前完成资格审核与学业成绩排名公示；中南大学材料学院要求成绩审核以截止 9 月 10 日 17:00 教务系统公布的加权平均分和专业排名为准。这意味着七八月是流程最密集的阶段。

因此可行的策略是按赛事周期分层：
\begin{itemize}
  \item \textbf{春季赛（3—6 月）}：时间上最理想，出结果早，能在暑假前拿到证书；
  \item \textbf{暑期赛（7—8 月）}：与夏令营冲突风险高，若必须参加，先确认赛程是否与目标院校夏令营时间重叠；
  \item \textbf{秋季赛（9—11 月）}：对当年推免已来不及，只能作为下一年的储备。
\end{itemize}

\textbf{加分口径上要现实一点}：安徽师范大学规定专长成绩最高不超过 10 分；中南大学材料学院规定各类型加分累计不超过 10 分，其中学院评价加分不超过 2 分。所以一次竞赛的收益是有限的，不要为此牺牲掉夏令营。

\textbf{判断方法}：先查本校的加分截止时点（多数在 8 月底前完成公示），再倒推报名时间。
\nref{S0540}
\nref{S0273}
\nref{S0646}
\nref{S0079}
\end{qa}

\iin{湖南大学}
\ith{竞赛认定}
\ist{背景提升}
\section{绩点提升}
\begin{qa}{M0181}{大二绩点只有3.3，还有机会靠大三拉起来吗}
\textbf{结论：有机会，但关键变量不是绝对分数，而是排名和“还剩多少门课能影响排名”。}

\textbf{第一件事：把问题换成排名问题。} 各校的门槛都是排名口径——内蒙古大学要求平均学分绩点排名不低于本专业前 50\%；安徽师范大学要求推免课程加权平均成绩排名前 35\%、前三年综合测评排名前 35\%；成都大学商学院要求本专业必修课程加权平均分排名前 20\%。GPA 3.3 在不同专业可能是前 10\%，也可能是后 50\%，只报绝对分数没有意义。

\textbf{第二件事：算清大三的学分权重。} 综合成绩中学业成绩的占比是刚性的：甘肃农业大学规定课程平均成绩占比不得低于 70\%，大连理工大学商学院为 85\%，成都大学商学院为 90\%。但\textbf{能改变的部分取决于大三课程占总计分课程的学分比例}——如果大三专业课学分比重大，拉动的空间就大；如果大部分课程在前两年已修完，提升幅度就有限。

\textbf{第三件事：注意计分口径。} 成都大学商学院规定课程成绩采用\textbf{首次获得的成绩}；中南大学材料学院规定补考和重修合格成绩按 60 分计算，且有不及格课程者不得推荐——说明“大三补救”要靠新修课程，而不是重修刷分。可行路径是先把握大三专业课考到能拿的最高分，再考虑专长加分。
\nref{S0544}
\nref{S0646}
\nref{S0079}
\nref{S0059}
\nref{S0549}
\nref{S0273}
\end{qa}

\iin{内蒙古大学}
\ith{绩点提升}
\ist{背景提升}
\begin{qa}{M0182}{怎么把绩点刷上去？选水课算不算投机取巧}
\textbf{结论：选给分高的课本身不违规，但“靠重修刷分”这条路在多数学校被明确堵死。} 规则重点治理的是重复计分，而不是选课偏好。

\textbf{刷分被堵死的具体条款}：
\begin{itemize}
  \item 成都大学商学院：课程成绩计算采用\textbf{首次获得的成绩}；
  \item 中南大学材料学院：\textbf{补考和重修合格的课程成绩按 60 分计算}，且推免生必须完成专业培养方案前 3 年所有教学环节且成绩合格，有不及格课程者不得推荐；
  \item 大连理工大学商学院：所有参评课程的成绩均以\textbf{一考成绩}计算，且一考不及格的必修课程不能超过 2 门；
  \item 甘肃农业大学：优先遴选推荐\textbf{无补考和重修记录}的学生，其次才遴选补考或重修课程不超过 1 门次的学生；
  \item 教育部《普通高等学校学生管理规定》第十八条：对通过补考、重修获得的成绩，应当予以标注。
\end{itemize}

\textbf{“选水课”这一侧的实际情况}：规则不禁止选修给分高的课程，但计分范围往往被限定——成都大学商学院只算必修课加权平均分，大连理工大学商学院按必修加专业选修模块计算，中南大学材料学院则用所有教学环节（含全校性选修课）加权平均裸分。真正的风险是把精力从必修课挪走导致排名下滑，以及留下一门不及格直接丧失资格。选课策略要在本校“推免课程清单”范围内做。
\nref{S0079}
\nref{S0273}
\nref{S0549}
\nref{S0059}
\nref{S0040}
\end{qa}

\iin{成都大学商学院}
\ith{绩点提升}
\ist{背景提升}
\begin{qa}{M0183}{为了保研选水课拉高绩点，会不会被学院查}
\textbf{结论：选课本身不会被查，因为规则不禁止；真正会被查、且后果严重的是材料造假。} 这两件事经常被混为一谈。

\textbf{学院在审查什么}：各校设置审核程序的目的是核验\textbf{加分项的真实性}，不是审查你的选课偏好。内蒙古大学要求学院成立专家审核小组（副教授以上职称不少于 5 人），对科研创新成果、论文、竞赛获奖进行审核鉴定，排除抄袭、造假、冒名及有名无实，组织公开答辩并录音录像；成都大学商学院由学院学术分委员会审核鉴定并组织公开答辩，答辩全程录音录像；中南大学材料学院、安徽师范大学均有同类规定。

\textbf{造假的实际后果}：甘肃农业大学规定，在推免过程中弄虚作假、有考试作弊或学术不端行为的学生，一经查实取消推免资格、录取资格和学籍，并记入《国家教育考试考生诚信档案》、通报所在学校、记入学生人事档案；安徽师范大学规定弄虚作假者一经查实将扣除该学院本年度相应的推免指标。

\textbf{程序上的救济渠道也存在}：湖南大学设复议申请表，南开大学设推免研究生复议申请表，学生可对推荐结果提出异议。真正需要担心的不是“选了水课”，而是选课导致必修课没修够或出现不及格，以及在成果材料上有任何不实描述。
\nref{S0544}
\nref{S0079}
\nref{S0273}
\nref{S0646}
\nref{S0059}
\nref{S0540}
\nref{S0561}
\end{qa}

\iin{内蒙古大学}
\ith{绩点提升}
\ist{背景提升}
\begin{qa}{M0184}{专业课给分低，是不是应该多修通识课拉平均}
\textbf{结论：不一定——各校对“哪些课计入排名”的规定差异极大，通识课在有些学校完全不计入，在另一些学校则计入。}

三种典型口径：

\par\vspace{2pt}\nobreak\begin{center}\footnotesize
\begin{tabularx}{\textwidth}{>{\raggedright\arraybackslash}X>{\raggedright\arraybackslash}X}
\toprule
\textbf{学校} & \textbf{计分课程范围} \\ \midrule
成都大学商学院 & 只算本专业培养方案中所有\textbf{必修课程}的加权平均分，通识选修课不计入 \\
大连理工大学商学院 & 必修课程 + 专业选修模块（需完成规定学分的 80\% 以上），并按一考成绩计算 \\
中南大学材料学院 & 在校学习期间\textbf{所有教学环节（含全校性选修课）}的加权平均裸分 \\
\bottomrule
\end{tabularx}
\end{center}\vspace{2pt}\par

\textbf{这个差异决定了策略完全不同}：按成都大学商学院的口径，多修通识课对排名毫无帮助；按中南大学的口径，全校性选修课确实计入，选给分高的课能产生实际效果。但中南大学另有条件——全校性选修课已达 6 个学分且及格的，不影响推免资格，说明这类课程存在最低学分要求，单纯堆数量未必划算。

\textbf{判断方法}：找学院推免实施细则里的“推免课程清单”或“计分课程范围”，逐条核对。大连理工大学商学院就明确列出不纳入计算的情形——课程属性为必修的导论课不计入、截至公布排名前一日未出一考成绩的课程不纳入。另外，专业课给分低是结构性问题，务实解法是避开连续多门高难度专业课的学期，而不是用通识课去稀释。
\nref{S0079}
\nref{S0549}
\nref{S0273}
\nref{S0646}
\end{qa}

\iin{成都大学商学院}
\ith{绩点提升}
\ist{背景提升}
\begin{qa}{M0222}{转专业后要补课，绩点被拉低怎么办}
\textbf{结论：转专业学生的绩点处理规则由各校自定，常见做法是“按现专业培养方案内的课程重新计算，不在范围内的课程不参与”。} 所以补课拉低的是你的绝对成绩，但不一定拉低排名。

\textbf{明确条款}：大连理工大学商学院规定，\textbf{“转专业学生按目前所学专业对应的必修课程范围计算，若所修课程不在规定课程范围内，则该课程不参与计算。”} 这条规则的实质是：原专业的课程成绩不计入，对新专业来说是“清零重算”。

\textbf{这意味着两种结果}：原专业成绩一般的学生重新起算反而占优；原专业成绩优异的学生则失去这部分优势。教育部《普通高等学校学生管理规定》第二十一条规定，学生可以申请转专业，学校应当制定具体办法并建立公平公正的标准、程序与公示制度——学分认定与成绩处理的认定权在学校。

\textbf{补救路径}：1）尽早向教务处索取“转专业学生课程替代与学分认定”的书面规则，确认哪些课能替代、如何计入；2）用剩余学期把现专业核心必修课成绩做上去，因为计分范围是现在的专业；3）补课时间挤压时优先保必修课，把通识与选修放到后面。安徽师范大学对交流生的要求可作参照——须通过课程替代、学分认定完成交流期间本专业必修课程，才可参加成绩计算和排名。
\nref{S0549}
\nref{S0040}
\nref{S0646}
\end{qa}

\iin{大连理工大学}
\ith{绩点提升}
\ist{背景提升}
\begin{qa}{M0235}{大二成绩下滑一门课会不会影响整体印象}
\textbf{结论：一门课的成绩下滑，影响主要是它在加权平均分里的权重，而不是评审对你“整体印象”的判断。} 但有一个例外情形需要特别小心。

\textbf{为什么影响通常有限}：推免排名是按加权平均分或平均学分绩点算的算术结果。大连理工大学商学院的计算方式是按课程学分加权，并规定所有参评课程成绩均以一考成绩计算、一考不及格的必修课不能超过 2 门；成都大学商学院要求本专业必修课程加权平均分排名前 20\%。一门 2—3 学分的课程，对总加权分的影响是百分之几的量级。

\textbf{真正的例外是“不及格”}：中南大学材料学院规定推免生必须完成专业培养方案前 3 年所有教学环节且成绩合格，\textbf{有不及格课程者不得推荐}；安徽师范大学要求所有推免课程无补考或重修；大连理工大学商学院对一考不及格必修课设了 2 门的上限。\textbf{一门课挂科的风险不是“印象差”，而是资格问题。}

\textbf{至于“印象”}：教育部《关于加强硕士研究生招生复试工作的指导意见》把“大学阶段学习情况及成绩”列为专业素质考核第一项，但看的是整体成绩而非单门课。正确应对是先确认这门课是否必修、是否及格、是否在计分范围内；若只是分数低，把后面的课做好即可，不必在个人陈述里主动解释。
\nref{S0549}
\nref{S0079}
\nref{S0273}
\nref{S0646}
\nref{S0046}
\end{qa}

\iin{大连理工大学}
\ith{绩点提升}
\ist{背景提升}
\section{论文产出}
\begin{qa}{M0192}{本科生发论文现实吗？要多久才能出一篇}
\textbf{结论：现实，但周期比多数人预期长，且学校只认“符合目录要求 + 署名位次达标”的那一部分。} 发一篇能真正用于保研的论文，从立项到录用通常需要一年以上。

\textbf{先看“能用”的标准有多严}：
\begin{itemize}
  \item 甘肃农业大学：须在高水平学术期刊（SCI、SSCI、EI、CSCD、CSSCI、《中文核心期刊要目总览》收录期刊或本科院校学报）以第一作者发表，或与指导教师联合发表（限指导教师和学生分别为第一、二作者），且\textbf{本校为第一署名单位}；
  \item 安徽师范大学：限独立作者或第一作者，且发表在学校“论文分类导向性期刊目录”\textbf{C2 及以上}期刊；
  \item 大连理工大学商学院：须以独立作者或第 1 作者身份发表或录用，且论文以\textbf{接收日期}为准，所有成果需本校为第一单位。
\end{itemize}

\textbf{再看它能换来多少分}：大连理工大学商学院的科研论文计分是 TOP 类期刊 30 分/篇、A 类 15 分、B 类 10 分、C 类 7 分、EI 检索期刊及中文核心 5 分、EI 检索会议论文 3 分，但\textbf{最多限 2 项}，且代表性成果相加超过 30 分按满分计；内蒙古大学的加分上限是学分绩点加分最高不超过 0.15。

\textbf{结论}：本科生发表论文是可行的，但它是一条“长周期、低分值的加分路”，而不是捷径。
\nref{S0059}
\nref{S0646}
\nref{S0549}
\nref{S0544}
\end{qa}

\iin{甘肃农业大学}
\ith{论文产出}
\ist{背景提升}
\begin{qa}{M0193}{发一篇SCI二作和一篇中文核心一作哪个更有用}
\textbf{结论：在多数学校的规则下，中文核心一作往往比 SCI 二作更“能用”，因为非导师一作的第二作者身份常常直接落在认定范围之外。}

\textbf{为什么二作会失效}：各校对论文的署名要求写得很死。甘肃农业大学规定须以第一作者发表，或与指导教师联合发表（限指导教师和学生分别为第一、二作者），且本校为第一署名单位；成都大学商学院只认定学生本人为第一作者；安徽师范大学限独立作者或第一作者。也就是说，\textbf{如果第二作者不是“导师一作、你二作”这种特定组合，这篇论文在加分认定中可能完全不算数}。

\textbf{分值上的对比}：中南大学材料学院的算法是——JCR 1 区 0.6 分、JCR 2 区 0.5 分、JCR 3/4 区及中文 A 类 0.4 分、EI 及中文 B 类 0.3 分；存在共同第一作者时，排名第一的按 80\% 计，其他共同一作按 20\% 计；\textbf{有多篇论文者不累计，取最高分}。大连理工大学商学院的论文项为 TOP 30 分、A 类 15 分、B 类 10 分、C 类 7 分、EI 及中文核心 5 分，最多限 2 项。

\textbf{所以判断顺序应该是}：先看署名位次是否达标（这是资格），再看期刊等级（这是分值）。一篇署名位次不被认可的高分区论文，收益是零；一篇中文核心一作，至少能拿到基础分。

\textbf{例外情形}：如果你的目标是申请海外或顶尖院系的直博，SCI 二作在导师推荐信和面试中仍有说服力——但那属于“申请材料价值”，不是“本校加分价值”。
\nref{S0059}
\nref{S0079}
\nref{S0646}
\nref{S0273}
\nref{S0549}
\end{qa}

\iin{甘肃农业大学}
\ith{论文产出}
\ist{背景提升}
\begin{qa}{M0194}{论文还在投没有录用，申请时怎么写在简历上}
\textbf{结论：可以写在简历和材料里，但必须如实标注状态；加分环节只认“已发表/已录用”，在投稿件通常不计分。}

\textbf{官方是允许标注状态的}：上海交通大学人工智能学院在夏令营报名说明中明确要求“论文请详细写明作者排序、录用状态（\textbf{正在投稿/已录用}）”——这说明在正式申请材料里写“在投”是被规则接受的表述方式。

\textbf{但加分口径要严格区分}：大连理工大学商学院规定论文以\textbf{接收日期}为准，专利以授权日期为准；西安电子科技大学计算机学院、光电工程学院在申请材料清单中写的是“发表或\textbf{已录用}的、能体现自身学术水平的学术论文”。两份文件合起来指向同一个口径：\textbf{加分认定看的是录用或发表，不看投稿。}

\textbf{所以写法建议}：
\begin{enumerate}
  \item 分栏列出——已发表/已录用 与 在投（Under Review）分开写；
  \item 在投项写清期刊/会议名称、投稿时间、当前状态（Submitted / Under Review / Major Revision）、你的作者位次；
  \item 附上投稿系统截图或稿件首页作为佐证，不要只写标题；
  \item 简历里避免用“发表”这个词描述在投稿件，用“投稿中”或“在审”。
\end{enumerate}

\textbf{面试中的加分点}：能说清你的投稿进度、审稿意见和你的修改方案，往往比“我有一篇论文”更能体现科研训练。
\nref{S0514}
\nref{S0549}
\nref{S0121}
\end{qa}

\iin{上海交通大学}
\ith{论文产出}
\ist{背景提升}
\begin{qa}{M0195}{保研申请时论文挂预印本算不算成果}
\textbf{结论：预印本在检索到的各校推免加分条款中未见被认可，它不算加分成果，但可以作为申请材料中的展示内容。}

\textbf{为什么加不了分}：各校认定的口径都指向正式发表或录用，并且要能对上期刊目录。甘肃农业大学列出的认可范围是 SCI、SSCI、EI、CSCD、CSSCI、《中文核心期刊要目总览》收录期刊或本科院校学报，并要求本校为第一署名单位；安徽师范大学要求发表在“论文分类导向性期刊目录”C2 及以上期刊；大连理工大学商学院规定论文以接收日期为准。预印本平台不属于上述任何一个目录。

\textbf{但它仍然有用}：一是可以作为“原创性工作成果”写进申请材料，各校夏令营材料普遍接受这一类；二是上海交通大学人工智能学院允许在报名中标注论文的“正在投稿”状态，说明未正式发表的稿件有合法展示位置；三是在面试里，一篇有完整方法、数据和结论的预印本能支撑一段高质量的科研对话。

\textbf{写法建议}：如实标注为“预印本 / Preprint”并给出链接或 DOI，不要写成“已发表”。同时准备好回答“为什么先挂预印本”“投稿进展如何”这类追问。

\textbf{风险提示}：部分期刊对预印本有政策限制，投稿前先确认目标期刊是否允许。
\nref{S0059}
\nref{S0646}
\nref{S0549}
\nref{S0514}
\end{qa}

\iin{甘肃农业大学}
\ith{论文产出}
\ist{背景提升}
\begin{qa}{M0223}{医学本科五年，前四年成绩怎么排，实习期怎么提升}
\textbf{结论：五年制医学专业通常用前四学年成绩排名，实习期（第五年）的核心任务不是刷绩点，而是产出科研与专长成果。}

\textbf{成绩口径是明确的}：多校文件都写了这条括号说明——甘肃农业大学要求“前三学年（五年制专业为前四学年）纳入课程平均成绩计算范围的课程考核合格”；成都大学商学院要求修读并获得“前三学年（五年制为前四学年）全部应修学分”；安徽师范大学的推荐对象为“2023 级（四年制）或 2022 级（五年制）”全日制普通本科毕业生，推免课程范围同样按五年制前四学年处理；中国科学院信息工程研究所在夏令营申请材料中也要求“大学本科前三年所修课程成绩单（五年制的提供前四年课程成绩单）”。

\textbf{实习期真正能提升的是两件事}：一是科研与专长加分——论文、竞赛、专利等加分项不受实习限制，五年制比四年制多出整整一年准备窗口；二是方向匹配度——临床医学、口腔医学等专业学位的接收往往要求相应专业的应届本科毕业生，实习经历在面试中直接相关。若第五年课程不在计分范围内，为刷分投入的精力性价比很低，具体以本校“推免课程清单”为准。

\textbf{医学类接收端的英语门槛通常更高}：浙江大学医学院规定申请直博需大学英语六级 480 分及以上，或新托福 90、雅思 6.0 及以上，且不认可家考成绩、拼分成绩。
\nref{S0059}
\nref{S0079}
\nref{S0646}
\nref{S0029}
\nref{S0219}
\end{qa}

\iin{甘肃农业大学}
\ith{医学学制}
\ist{背景提升}

\section{科研过程与成果呈现}
\begin{qa}{U202610210401}{本科科研项目尚未结题，申请材料里怎样写阶段性成果？}
可以写，但要把“在研”写在标题或首句，并给出材料截止日，不能把计划、预期结果写成已完成。推荐四段式：①研究问题与项目背景（一两句）；②你本人承担的具体任务（数据清洗、实验、访谈、代码模块等，用动词写清）；③截至材料提交日已完成的步骤和可核验产物（阶段报告、代码仓库、实验记录、问卷回收数、中期检查表）；④尚未完成的部分、原因和下一节点。

产物要能经得起追问：准备一段一分钟说明“我做了什么、遇到什么问题、怎么解决、目前证据在哪”。涉及导师课题、未公开数据或企业保密内容，先取得负责人同意并脱敏，不把团队成果写成个人独立完成；也不要用“预计投稿”“拟申请专利”替代现有结果。若项目中途终止，如实说明终止原因、你负责部分的完成度和可交接材料，比强行包装成结题更有说服力。
\qscope{材料写法建议；项目保密审查和成果署名以项目负责人及学校规定为准。}
\end{qa}
\begin{qa}{U202610210402}{项目代码或数据暂时不能公开，面试时怎样证明自己的实际贡献？}
先确认项目负责人和学校允许展示的范围，保密要求优先于申请展示。证明贡献不要靠“我做了很多”这种描述，准备一套可核验的四件套：一，任务边界，用一页纸写清项目目标、你负责的模块、输入输出和与其他成员的接口；二，过程证据，提交记录、实验日志、版本号、issue 或评审意见（去掉敏感信息）；三，结果证据，可公开的指标、对照实验、用户反馈或导师签字的任务说明；四，替代样例，用去标识化的小数据或不含业务逻辑的代码骨架现场讲解。面试前把材料给导师或项目负责人审阅，遮去姓名、账号、患者或受试者信息、未公开数据和合作方资料；不能展示原件时直接说明限制，用方法、决策过程和可公开样例回答追问，不伪造截图、不夸大署名，也不要把“团队成果”说成个人独立完成。
\qscope{通用材料准备建议；保密义务、数据出境和知识产权归属以项目协议及学校规定为准。}
\end{qa}
\begin{qa}{U202610210403}{还没有论文或获奖，怎么把课程和科研训练整理成可信的申请证据？}
不要为了填简历空白临时购买论文、挂名成果或把课程作业包装成已发表研究。可以选一项真正做过的课程项目、实验训练或独立阅读，整理成“问题—方法—本人任务—结果—局限”一页记录，并保留报告、代码、实验日志或教师反馈等原始材料。面试重点是能解释为什么这样设计、遇到什么问题、结论有哪些边界；没有正式成果就明确写“课程项目”或“科研训练”，不写成论文、课题负责人或已结题项目。是否计入综合评价仍以本校和目标院系当年材料规则为准。
\end{qa}
